% GENERATED FILE. Edit canonical lessons, then run npm run generate:books.
% canonical-source-hash: fnv1a64:8d5ccc1c28e810cb
% canonical-lessons: ML-C481-risorttu, ML-C481-satram, ML-C481-tamasasthalam, ML-C481-tamasakkaran, ML-C481-risapsan

\chapter{Finding a Room}
\label{ch:ml-finding-a-room}
\hlchaptermodality{481}

\begin{chapteropening}
\textbf{By the end of this chapter:} \emph{I can say a resort, an inn, accommodation, a resident and a reception desk.}

Complete the last lesson of chapter 481: I can say a resort, an inn, accommodation, a resident and a reception desk.
\end{chapteropening}

\section[\texorpdfstring{\ml{റിസോർട്ട്} (ṟisōrṭṭŭ)}{റിസോർട്ട് (ṟisōrṭṭŭ)}]{\ml{റിസോർട്ട്} (ṟisōrṭṭŭ) — a resort}

\label{lesson:ML-C481-risorttu}



\begin{quote}
Before the new one: say the Malayalam for a locker, then the Malayalam for to keep aside.
\end{quote}

\subsection*{You'll want to know: \ml{റിസോർട്ട്}}
\textbf{\ml{റിസോർട്ട്}} — \emph{ṟisōrṭṭŭ} — \textquotedblleft{}a resort\textquotedblright{}.

\begin{quote}
\ml{കായലിന്റെ} \ml{തീരത്ത്} \ml{ഒരു} \ml{പുതിയ} \ml{റിസോർട്ട്} \ml{ഉണ്ട്}.
\end{quote}

\begin{quote}
\emph{kāyalinṟe tīrattŭ oru putiya ṟisōrṭṭŭ uṇṭŭ.}
\end{quote}

\begin{quote}
There is a new resort on the shore of the backwaters.
\end{quote}

A resort is bigger and dearer than an ordinary \ml{ഹോട്ടൽ}.

\begin{cousinweb}
\textbf{\ml{റിസോർട്ട്}} (\emph{ṟisōrṭṭŭ}) was borrowed from English \emph{resort}, a place where people stay on holiday.
\end{cousinweb}

\begin{grammarlens}[title={What you've built}]
One more word for getting around.
\end{grammarlens}

\subsection*{Your turn}
\begin{itemize}
\raggedright
  \item \emph{Say it:} \emph{ṟisōrṭṭŭ}
  \item \emph{Say it:} \emph{ṟisōrṭṭŭ}, once more
  \item \emph{Say it:} \emph{kāyalinṟe tīrattŭ oru putiya ṟisōrṭṭŭ uṇṭŭ.} — There is a new resort on the shore of the backwaters.
  \item \emph{Recall:} say the Malayalam for a locker, then the Malayalam for to keep aside, then say \emph{ṟisōrṭṭŭ} again
\end{itemize}

\begin{tcolorbox}[breakable,colback=teal!4,colframe=teal!35!black,title={Before you move on}]
What does \ml{റിസോർട്ട്} mean? (\textquotedblleft{}A resort\textquotedblright{}.) Say it once more.
\end{tcolorbox}
\section[\texorpdfstring{\ml{സത്രം} (satraṁ)}{സത്രം (satraṁ)}]{\ml{സത്രം} (satraṁ) — an inn; a rest house}

\label{lesson:ML-C481-satram}



\begin{quote}
Before the new one: say the Malayalam for to keep aside, then the Malayalam for a resort.
\end{quote}

\subsection*{You'll want to know: \ml{സത്രം}}
\textbf{\ml{സത്രം}} — \emph{satraṁ} — \textquotedblleft{}an inn; a rest house\textquotedblright{}.

\begin{quote}
\ml{പണ്ട്} \ml{യാത്രക്കാർ} \ml{വഴിയിലെ} \ml{സത്രങ്ങളിൽ} \ml{തങ്ങിയിരുന്നു}.
\end{quote}

\begin{quote}
\emph{paṇṭŭ yātrakkār vaḻiyile satraṅṅaḷil taṅṅiyirunnu.}
\end{quote}

\begin{quote}
Long ago travellers used to stay at inns along the way.
\end{quote}

Plural: \ml{സത്രങ്ങൾ}, \emph{satraṅṅaḷ}.

A -\ml{ം} noun: the plural is -\ml{ങ്ങൾ}, \ml{സത്രങ്ങൾ}.

\begin{cousinweb}
\textbf{\ml{സത്രം}} (\emph{satraṁ}) is borrowed from Sanskrit \emph{sattra}, first a long ritual session and later also a house or shelter (Monier-Williams).
\end{cousinweb}

\begin{grammarlens}[title={What you've built}]
One more word for getting around.
\end{grammarlens}

\subsection*{Your turn}
\begin{itemize}
\raggedright
  \item \emph{Say it:} \emph{satraṁ}
  \item \emph{Say it:} \emph{satraṁ}, once more
  \item \emph{Say it:} \emph{paṇṭŭ yātrakkār vaḻiyile satraṅṅaḷil taṅṅiyirunnu.} — Long ago travellers used to stay at inns along the way.
  \item \emph{Recall:} say the Malayalam for to keep aside, then the Malayalam for a resort, then say \emph{satraṁ} again
\end{itemize}

\begin{tcolorbox}[breakable,colback=teal!4,colframe=teal!35!black,title={Before you move on}]
What does \ml{സത്രം} mean? (\textquotedblleft{}An inn; a rest house\textquotedblright{}.) Say it once more.
\end{tcolorbox}
\section[\texorpdfstring{\ml{താമസസ്ഥലം} (tāmasasthalaṁ)}{താമസസ്ഥലം (tāmasasthalaṁ)}]{\ml{താമസസ്ഥലം} (tāmasasthalaṁ) — accommodation; a place to stay}

\label{lesson:ML-C481-tamasasthalam}



\begin{quote}
Before the new one: say the Malayalam for a resort, then the Malayalam for an inn.
\end{quote}

\subsection*{You'll want to know: \ml{താമസസ്ഥലം}}
\textbf{\ml{താമസസ്ഥലം}} — \emph{tāmasasthalaṁ} — \textquotedblleft{}accommodation; a place to stay\textquotedblright{}.

\begin{quote}
\ml{നഗരത്തിൽ} \ml{വില} \ml{കുറഞ്ഞ} \ml{താമസസ്ഥലം} \ml{കിട്ടാൻ} \ml{പ്രയാസമാണ്}.
\end{quote}

\begin{quote}
\emph{nagarattil vila kuṟañña tāmasasthalaṁ kiṭṭān prayāsamāṇŭ.}
\end{quote}

\begin{quote}
Cheap accommodation is hard to find in the city.
\end{quote}

\ml{താമസം}, a stay, plus \ml{സ്ഥലം}, a place.

\begin{cousinweb}
\textbf{\ml{താമസസ്ഥലം}} (\emph{tāmasasthalaṁ}) is a Malayalam compound built from \ml{താമസം}, which means both a delay and a stay and is usually traced to Sanskrit \emph{tāmasa}, \textquotedblleft{}dark, inert\textquotedblright{} (CDIAL 5775), and \ml{സ്ഥലം}, a place, from Sanskrit \emph{sthala} (Monier-Williams).
\end{cousinweb}

\begin{grammarlens}[title={What you've built}]
One more word for getting around.
\end{grammarlens}

\subsection*{Your turn}
\begin{itemize}
\raggedright
  \item \emph{Say it:} \emph{tāmasasthalaṁ}
  \item \emph{Say it:} \emph{tāmasasthalaṁ}, once more
  \item \emph{Say it:} \emph{nagarattil vila kuṟañña tāmasasthalaṁ kiṭṭān prayāsamāṇŭ.} — Cheap accommodation is hard to find in the city.
  \item \emph{Recall:} say the Malayalam for a resort, then the Malayalam for an inn, then say \emph{tāmasasthalaṁ} again
\end{itemize}

\begin{tcolorbox}[breakable,colback=teal!4,colframe=teal!35!black,title={Before you move on}]
What does \ml{താമസസ്ഥലം} mean? (\textquotedblleft{}Accommodation; a place to stay\textquotedblright{}.) Say it once more.
\end{tcolorbox}
\section[\texorpdfstring{\ml{താമസക്കാരൻ} (tāmasakkāran)}{താമസക്കാരൻ (tāmasakkāran)}]{\ml{താമസക്കാരൻ} (tāmasakkāran) — a resident; a lodger}

\label{lesson:ML-C481-tamasakkaran}



\begin{quote}
Before the new one: say the Malayalam for an inn, then the Malayalam for accommodation.
\end{quote}

\subsection*{You'll want to know: \ml{താമസക്കാരൻ}}
\textbf{\ml{താമസക്കാരൻ}} — \emph{tāmasakkāran} — \textquotedblleft{}a resident; a lodger\textquotedblright{}.

\begin{quote}
\ml{ആ} \ml{വീട്ടിലെ} \ml{പുതിയ} \ml{താമസക്കാരൻ} \ml{ഒരു} \ml{അധ്യാപകനാണ്}.
\end{quote}

\begin{quote}
\emph{ā vīṭṭile putiya tāmasakkāran oru adhyāpakanāṇŭ.}
\end{quote}

\begin{quote}
The new resident of that house is a teacher.
\end{quote}

Plural: \ml{താമസക്കാർ}, \emph{tāmasakkār}.

The ending -\ml{കാരൻ} means one who: \ml{താമസം}, a stay, so one who stays. The plural is \ml{താമസക്കാർ}.

\begin{cousinweb}
\textbf{\ml{താമസക്കാരൻ}} (\emph{tāmasakkāran}) is built from \ml{താമസം}, a stay, usually traced to Sanskrit \emph{tāmasa}, \textquotedblleft{}dark, inert\textquotedblright{} (CDIAL 5775), with -\ml{കാരൻ}, the Malayalam ending for a person who does something.
\end{cousinweb}

\begin{grammarlens}[title={What you've built}]
One more word for getting around.
\end{grammarlens}

\subsection*{Your turn}
\begin{itemize}
\raggedright
  \item \emph{Say it:} \emph{tāmasakkāran}
  \item \emph{Say it:} \emph{tāmasakkāran}, once more
  \item \emph{Say it:} \emph{ā vīṭṭile putiya tāmasakkāran oru adhyāpakanāṇŭ.} — The new resident of that house is a teacher.
  \item \emph{Recall:} say the Malayalam for an inn, then the Malayalam for accommodation, then say \emph{tāmasakkāran} again
\end{itemize}

\begin{tcolorbox}[breakable,colback=teal!4,colframe=teal!35!black,title={Before you move on}]
What does \ml{താമസക്കാരൻ} mean? (\textquotedblleft{}A resident; a lodger\textquotedblright{}.) Say it once more.
\end{tcolorbox}
\section[\texorpdfstring{\ml{റിസപ്ഷൻ} (ṟisapṣan)}{റിസപ്ഷൻ (ṟisapṣan)}]{\ml{റിസപ്ഷൻ} (ṟisapṣan) — a reception desk}

\label{lesson:ML-C481-risapsan}



\begin{quote}
Before the new one: say the Malayalam for accommodation, then the Malayalam for a resident.
\end{quote}

\subsection*{You'll want to know: \ml{റിസപ്ഷൻ}}
\textbf{\ml{റിസപ്ഷൻ}} — \emph{ṟisapṣan} — \textquotedblleft{}a reception desk\textquotedblright{}.

\begin{quote}
\ml{ഹോട്ടലിലെ} \ml{റിസപ്ഷനിൽ} \ml{ഞാൻ} \ml{എന്റെ} \ml{പേര്} \ml{പറഞ്ഞു}.
\end{quote}

\begin{quote}
\emph{hōṭṭalile ṟisapṣanil ñān enṟe pērŭ paṟaññu.}
\end{quote}

\begin{quote}
I gave my name at the hotel reception.
\end{quote}

A wedding reception is also a \ml{റിസപ്ഷൻ}.

\begin{cousinweb}
\textbf{\ml{റിസപ്ഷൻ}} (\emph{ṟisapṣan}) was borrowed from English \emph{reception}, the front desk of a hotel or office.
\end{cousinweb}

\begin{grammarlens}[title={What you've built}]
One more word for getting around.
\end{grammarlens}

\subsection*{Your turn}
\begin{itemize}
\raggedright
  \item \emph{Say it:} \emph{ṟisapṣan}
  \item \emph{Say it:} \emph{ṟisapṣan}, once more
  \item \emph{Say it:} \emph{hōṭṭalile ṟisapṣanil ñān enṟe pērŭ paṟaññu.} — I gave my name at the hotel reception.
  \item \emph{Recall:} say the Malayalam for accommodation, then the Malayalam for a resident, then say \emph{ṟisapṣan} again
\end{itemize}

\begin{tcolorbox}[breakable,colback=teal!4,colframe=teal!35!black,title={Before you move on}]
What does \ml{റിസപ്ഷൻ} mean? (\textquotedblleft{}A reception desk\textquotedblright{}.) Say it once more.
\end{tcolorbox}
